\section{Canonical tensor-product coefficient и threshold saturation}

Project-wide archaeology оставила mixed relative curvature как последний
внутренний кандидат. Однако до проверки новых contractions необходимо
зафиксировать коэффициент самого простого product-triple механизма, поскольку
его нельзя считать свободным parameter.

\subsection{Канонический graded product}

Для произведения чётной spectral triple с auxiliary factor стандартный
Dirac operator имеет вид
\begin{equation}
 D_\times=D_1\otimes I+\Gamma_1\otimes D_2,
 \qquad \{D_1,\Gamma_1\}=0.
\end{equation}
Отсюда без дополнительного предположения следует
\begin{equation}
 D_\times^2=D_1^2\otimes I+I\otimes D_2^2.
\end{equation}
Поэтому mixed часть quartic spectral trace равна
\begin{equation}
 \frac12\Tr D_\times^4\big|_{\rm mixed}
 =\Tr(D_1^2)\Tr(D_2^2).
 \label{eq:ps-product-mixed-trace}
\end{equation}
Именно эта product prescription используется в стандартных построениях
real spectral triples; возможные KO-sign conventions меняют real structure,
но не тождество для квадрата данного graded Dirac operator.

Пусть оба factors заданы odd self-adjoint completions комплексных blocks
\(M\) и \(X\). Тогда
\begin{equation}
 \Tr D_1^2=2\|M\|_F^2,
 \qquad
 \Tr D_2^2=2\|X\|_F^2,
\end{equation}
и \eqref{eq:ps-product-mixed-trace} принимает вид
\begin{equation}
 \boxed{
 \frac12\Tr D_\times^4\big|_{\rm mixed}
 =4\|M\|_F^2\|X\|_F^2.}
 \label{eq:ps-canonical-tensor-c4}
\end{equation}
Следовательно, canonical tensor coefficient равен точно
\(c_{\rm can}=4\). Удвоение particle--conjugate sector в KO-dimension six
и последующий physical half-trace одновременно удваивают и делят trace,
поэтому оставляют тот же коэффициент. Один oriented self-adjoint curvature
block даёт ещё меньший coefficient \(c=2\).

\subsection{Сравнение с точным vacuum threshold}

Предыдущий Hessian gate установил, что для mixed portal
\begin{equation}
 V_{\rm mix}=c\,\Tr(M_R^\dagger M_R)\|Y\|_F^2
\end{equation}
на rank-one background
\(\Tr(M_R^\dagger M_R)=1/2\) induced scalar shift равен
\(\zeta=c/2\). Строгая устойчивость всех восьми project-
\(\phi\) directions требует
\begin{equation}
 \zeta>2
 \qquad\Longleftrightarrow\qquad
 c>4.
\end{equation}
Каноническое значение \(c=4\) только насыщает границу: четыре eigenvalues
становятся положительными, но ещё четыре остаются нулевыми. Таким образом,
full-rank tensor contraction устраняет отрицательные modes, но не создаёт
strict local minimum.

\begin{theorem}[Canonical tensor threshold saturation]
Канонический graded product двух odd finite blocks фиксирует mixed quartic
coefficient \(c=4\). После KO6 physical half-trace коэффициент не меняется.
Поэтому product spectral action достигает границы \(\zeta=2\), но не проходит
условие strict \(\phi\)-stability \(\zeta>2\).
\end{theorem}

\subsection{Почему multiplicity не является исправлением}

Для \(k\) идентичных auxiliary copies trace формально даёт
\(c_k=4k\). Уже \(k=2\) обеспечило бы \(c=8\), однако copy space имеет
commutant
\begin{equation}
 M_k(\mathbb C).
\end{equation}
При \(k>1\) он несводим не к скалярам, а содержит дополнительные matrix
endomorphisms. Это противоречит ранее выведенному one-copy irreducibility
selector и возвращает запрещённую multiplicity fork.

\begin{nogocriterion}[Identical-copy tensor rescue]
Нельзя поднимать canonical coefficient выше четырёх простым повторением
одинакового auxiliary carrier. Численный threshold тогда проходит только
ценой non-scalar copy commutant и нарушения irreducibility.
\end{nogocriterion}

\subsection{Оставшаяся развилка}

Tensor-product route остаётся открытым только в двух более узких формах:
\begin{enumerate}
 \item non-identical irreducible auxiliary factor, для которого joint
 commutant остаётся scalar, а induced coefficient превышает четыре;
 \item higher spectral moment с заранее выведенным coefficient и полной
 повторной radial stationarization, а не добавленный вручную quartic weight.
\end{enumerate}
Если ни одна ветвь не поднимает четыре нулевые modes без нового continuous
input, canonical tensor-product rescue закрывается.

Литературная product formula сверена с arXiv:math-ph/9902029 и
arXiv:1011.4456. Machine audit сохранён в
\path{s2t_v4_pati_salam_tensor_product_coefficient_gate.py} и
\path{s2t_v4_pati_salam_tensor_product_coefficient_gate_results.json}.